\documentclass[11pt,a4paper]{article}
\usepackage[utf8]{inputenc}
\usepackage[T1]{fontenc}
\usepackage[italian]{babel}
\usepackage[sc]{mathpazo}
\linespread{1.06}
\usepackage{amsmath,amssymb,amsthm,mathtools}
\usepackage[margin=2.5cm]{geometry}
\usepackage{booktabs,array,tabularx}
\usepackage[shortlabels]{enumitem}
\usepackage{microtype}
\usepackage[dvipsnames]{xcolor}
\usepackage[colorlinks=true,linkcolor=MidnightBlue,citecolor=MidnightBlue,urlcolor=MidnightBlue]{hyperref}
\usepackage[italian,capitalise,nameinlink]{cleveref}

\theoremstyle{plain}
\newtheorem{theorem}{Teorema}[section]
\newtheorem{lemma}[theorem]{Lemma}
\newtheorem{proposition}[theorem]{Proposizione}
\newtheorem{corollary}[theorem]{Corollario}
\newtheorem{conjecture}[theorem]{Congettura}
\theoremstyle{definition}
\newtheorem{definition}[theorem]{Definizione}
\newtheorem{hypothesis}[theorem]{Ipotesi}
\theoremstyle{remark}
\newtheorem{remark}[theorem]{Osservazione}
\crefname{theorem}{Teorema}{Teoremi}
\crefname{lemma}{Lemma}{Lemmi}
\crefname{proposition}{Proposizione}{Proposizioni}
\crefname{corollary}{Corollario}{Corollari}
\crefname{conjecture}{Congettura}{Congetture}
\crefname{definition}{Definizione}{Definizioni}
\crefname{hypothesis}{Ipotesi}{Ipotesi}
\crefname{remark}{Osservazione}{Osservazioni}
\crefname{section}{Sezione}{Sezioni}
\crefname{appendix}{Appendice}{Appendici}
\crefname{table}{Tabella}{Tabelle}
\crefname{equation}{}{}

\newcommand{\ket}[1]{\lvert #1\rangle}
\newcommand{\bra}[1]{\langle #1\rvert}
\newcommand{\braket}[2]{\langle #1\vert #2\rangle}
\newcommand{\Stab}{\mathrm{Stab}}
\newcommand{\Paulis}{\mathcal P}
\newcommand{\Fb}{\mathbb F}
\newcommand{\C}{\mathbb C}
\newcommand{\Q}{\mathbb Q}
\newcommand{\Z}{\mathbb Z}
\newcommand{\E}{\mathop{\mathbb E}}
\newcommand{\poly}{\mathrm{poly}}
\newcommand{\polylog}{\mathrm{polylog}}
\newcommand{\ChiP}{\chi_{\otimes}}
\newcommand{\Cl}{\mathcal C}
\newcommand{\eps}{\varepsilon}
\newcommand{\Tst}{\ket{T}}
\newcommand{\Hst}{\ket{H}}
\newcommand{\Gram}{\mathsf G}

\title{\textbf{Lower bound sul rango stabilizer di $\ket{T}^{\otimes n}$}\\[4pt]
\large Stato dell'arte, lower bound esponenziali in modelli ristretti e una congettura in termini di entropia stabilizer}
\author{Nota di lavoro\thanks{Bozza preparata con Claude Code. Le dimostrazioni delle Sezioni~\ref{sec:cond}--\ref{sec:sre} sono elementari ma vanno rilette con cura. Le affermazioni sulla letteratura sono state controllate sugli abstract arXiv e, dove indicato, sul testo; non sempre sulle dimostrazioni. Gli script delle verifiche numeriche (\cref{app:num}) sono disponibili a parte.}}
\date{26 novembre 2026}

\begin{document}
\maketitle

\begin{abstract}
\noindent Il rango stabilizer $\chi(\Tst^{\otimes n})$ governa il costo delle simulazioni classiche di circuiti Clifford$+T$ basate su decomposizioni stabilizer. Ci si aspetta che cresca esponenzialmente, ma il miglior lower bound noto è $\Omega(n^2/\log n)$, contro un upper bound $O(2^{0.3963\,n})$. In questa nota:
(i)~riassumiamo lo stato dell'arte, con particolare attenzione al ruolo dell'extent;
(ii)~osserviamo che un lower bound $2^{\Omega(n)}$ segue da una versione non uniforme di \#ETH;
(iii)~dimostriamo lower bound esponenziali \emph{incondizionati} in tre modelli ristretti (decomposizioni in stati prodotto, termini con molti qubit non entangled, decomposizioni ben condizionate) tramite un lemma di contrazione, e spieghiamo perché il metodo si ferma sui termini entangled;
(iv)~riformuliamo il problema tramite l'entropia di Rényi stabilizer $M_2$, che coincide con la norma di Gowers $U^3$ delle ampiezze, e formuliamo la congettura $M_2(\psi)\le C\log_2\chi(\psi)+D$. La congettura implica $\chi(\Tst^{\otimes n})=2^{\Omega(n)}$. A meno di costanti, è equivalente alla versione polinomiale della relazione rango--fidelity lasciata aperta da Sinha e Lovitz.
\end{abstract}

\tableofcontents

\section{Il problema}\label{sec:problem}

\subsection{Definizioni}
Sia $\Stab_n$ l'insieme degli stati stabilizer puri e normalizzati su $n$ qubit, e sia $\Paulis_n=\{I,X,Y,Z\}^{\otimes n}$. Per uno stato $\ket\psi$ definiamo:
\begin{itemize}[leftmargin=1.5em]
\item \textbf{rango stabilizer}: $\chi(\psi)=\min\bigl\{k:\ \ket\psi=\sum_{j=1}^k c_j\ket{s_j},\ s_j\in\Stab_n,\ c_j\in\C\bigr\}$;
\item \textbf{rango approssimato} (per $\|\psi\|=1$, $0\le\delta<1$): $\chi_\delta(\psi)=\min\{\chi(\phi):\ \|\psi-\phi\|\le\delta\}$;
\item \textbf{extent} \cite{BBCCGH19}: $\xi(\psi)=\min\bigl\{\|c\|_1^2:\ \ket\psi=\sum_j c_j\ket{s_j}\bigr\}$;
\item \textbf{fidelity stabilizer}: $F(\psi)=\max_{s\in\Stab_n}|\braket{s}{\psi}|^2$;
\item \textbf{entropia di Rényi stabilizer} \cite{LOH22}: $M_2(\psi)=-\log_2\Bigl(2^{-n}\sum_{P\in\Paulis_n}\bra\psi P\ket\psi^4\Bigr)$.
\end{itemize}
Lo stato magic è $\Tst=(\ket0+e^{i\pi/4}\ket1)/\sqrt2=T\ket+$. È Clifford-equivalente a
\begin{equation}\label{eq:H}
\Hst=\cos\tfrac\pi8\,\ket0+\sin\tfrac\pi8\,\ket1 .
\end{equation}
Tutte le quantità sopra sono invarianti per unitarie di Clifford. Quindi $\chi(\Tst^{\otimes n})=\chi(\Hst^{\otimes n})$, e lo stesso vale per le altre quantità; useremo l'uno o l'altro stato a seconda della comodità. Fidelity ed extent sono moltiplicativi sui prodotti di stati a un qubit \cite{BBCCGH19}, mentre $M_2$ è additiva. Si ha quindi
\begin{equation}\label{eq:values}
\begin{gathered}
F(\Hst^{\otimes n})=\cos^{2n}\tfrac\pi8,\qquad
\xi(\Hst^{\otimes n})=\cos^{-2n}\tfrac\pi8=(4-2\sqrt2)^n\approx 2^{0.2284\,n},\\
M_2(\Hst^{\otimes n})=n\log_2\tfrac43\approx 0.415\,n .
\end{gathered}
\end{equation}

\subsection{Relazioni elementari}
Valgono le seguenti relazioni.
\begin{enumerate}[label=(R\arabic*),leftmargin=2.5em]
\item\label{R1} $\chi_\delta(\psi)\le\chi(\psi)$. Inoltre $\chi_\delta(\psi)\le 1+\xi(\psi)/\delta^2$ per sparsificazione \cite{BBCCGH19}. Un lower bound su $\chi_\delta$ è quindi anche un lower bound su $\chi$.
\item\label{R2} (Dualità) Per ogni vettore $\omega$ vale $\xi(\psi)\ge|\braket{\omega}{\psi}|^2/\max_{s\in\Stab_n}|\braket{\omega}{s}|^2$. Con $\omega=\psi$ si ottiene $\xi(\psi)\ge 1/F(\psi)$.
\item\label{R3} (Monotonia) Se $\phi\notin\Stab_1$, allora $\chi(\phi^{\otimes(n+1)})\ge\chi(\phi^{\otimes n})$. Per vederlo, si contrae un qubit con uno stato $u\in\Stab_1$ tale che $\braket{u}{\phi}\neq0$: ogni termine diventa uno stato stabilizer oppure $0$.
\end{enumerate}

\subsection{Gate contro stati}
Il rango di Clifford di un operatore è $\chi_\Cl(U)=\min\{k:\ U=\sum_{j=1}^k c_jC_j,\ C_j\ \text{Clifford}\}$. Se $U=\sum_jc_jC_j$, lo stato di Choi $(U\otimes I)\ket{\Phi}$ è combinazione di $k$ stati stabilizer $(C_j\otimes I)\ket\Phi$.

Per il gate $T$ vale $(T\otimes I)\ket{\Phi^+}=(\ket{00}+e^{i\pi/4}\ket{11})/\sqrt2=\mathrm{CNOT}\,(\Tst\otimes\ket0)$. Quindi lo stato di Choi di $T^{\otimes n}$ è Clifford-equivalente a $\Tst^{\otimes n}\otimes\ket0^{\otimes n}$, e
\[
\chi_\Cl(T^{\otimes n})\ \ge\ \chi(\Tst^{\otimes n}).
\]
Ogni lower bound sugli stati vale dunque anche per la decomposizione del gate $T^{\otimes n}$. Nel seguito lavoriamo con gli stati.

\subsection{Perché interessa}
Data una decomposizione di $\Tst^{\otimes t}$ in $k$ termini, le ampiezze di un circuito Clifford$+T$ con $t$ gate $T$ si calcolano in tempo $k\cdot\poly(n,t)$ \cite{BG16,BBCCGH19}. Il rango stabilizer è quindi il parametro che controlla questi simulatori, e ci si aspetta $\chi(\Tst^{\otimes t})=2^{\Theta(t)}$. Il problema è dimostrarlo.

\section{Stato dell'arte}\label{sec:sota}

\begin{table}[ht]
\centering\small
\renewcommand{\arraystretch}{1.25}
\begin{tabularx}{\textwidth}{@{}p{4.3cm}Xp{3.0cm}@{}}
\toprule
Riferimento & Risultato & Tipo \\
\midrule
\multicolumn{3}{@{}l}{\emph{Upper bound}}\\
Bravyi--Smith--Smolin \cite{BSS16} & $\chi(\Tst^{\otimes 2})=2$, $\ \chi(\Tst^{\otimes 6})\le 7$ & esatto \\
Qassim--Pashayan--Gosset \cite{QPG21} & $\chi(\Tst^{\otimes n})=O(2^{\alpha n})$ con $\alpha\le 0.3963$ & esatto \\
Bravyi--Gosset; Bravyi et al.\ \cite{BG16,BBCCGH19} & $\chi_\delta\le 1+\xi/\delta^2=O(2^{0.2284\,n}/\delta^2)$ & approssimato \\
\midrule
\multicolumn{3}{@{}l}{\emph{Lower bound}}\\
Bravyi--Smith--Smolin \cite{BSS16} & $\Omega(\sqrt n)$ & esatto \\
Peleg--Shpilka--Volk \cite{PSV22} & $\chi\ge n/100$;\ \ $\chi_\delta=\Omega(\sqrt n/\log n)$ & esatto / appross. \\
Lovitz--Steffan \cite{LS22} & $\chi(\phi^{\otimes n})=\Omega(n/\log n)$, $\ \chi_\delta(\phi^{\otimes n})\ge \sqrt n/(2\log_2 n)$ per ogni $\phi\notin\Stab_1$ & esatto / appross. \\
Labib \cite{Lab22} & $\Omega(n)$ per l'analogo qudit, via analisi di Fourier di ordine superiore & esatto \\
Mehraban--Tahmasbi \cite{MT23} & $\chi_\delta(\Tst^{\otimes n})=\Omega(n^2)/\polylog(n)$ per $0<\delta<1$ & approssimato \\
Kalra--Sinha \cite{KS25} & $\chi(\psi)\le k\Rightarrow F(\psi)\ge 2^{-O(k\log k)}$, da cui $\Omega(n/\log n)$ & esatto \\
Sinha--Lovitz \cite{SL26} & $\chi(\psi)\le k\Rightarrow \xi(\psi)\le 2^{O(\sqrt{k\log k})}$;\ \ $\chi_\delta(\phi^{\otimes n})=\Omega(n^2/\log n)$ per ogni $\phi\notin\Stab_1$ & approssimato \\
\midrule
\multicolumn{3}{@{}l}{\emph{Valori esatti per $n$ piccolo}}\\
Labib--Russo \cite{LR26} & $\chi(\Tst^{\otimes4})=4$; per lo stato ``faccia'' $\chi=2,2,3,3$ per $m=1,\dots,4$ copie & ricerca al computer \\
\bottomrule
\end{tabularx}
\caption{Bound noti per il rango stabilizer delle potenze tensoriali di stati magic a un qubit. Per \ref{R1}, i lower bound sul rango approssimato valgono anche per il rango esatto.}
\label{tab:sota}
\end{table}

\paragraph{Il gap.} Il miglior lower bound noto, esatto o approssimato, è $\Omega(n^2/\log n)$ \cite{MT23,SL26}. Il miglior upper bound esatto è $O(2^{0.3963n})$ \cite{QPG21}, quello approssimato è $O(2^{0.2284n})$. Un lower bound superpolinomiale è aperto \cite{CC26,SL26}.

\paragraph{Da dove viene il quadratico.} Sinha e Lovitz dimostrano che uno stato di rango $k$ ha extent al più $2^{O(\sqrt{k\log k})}$. Combinato con $\xi(\Hst^{\otimes n})=2^{0.2284n}$ dà $\sqrt{k\log k}\gtrsim n$, cioè $k\gtrsim n^2/\log n$. Ogni miglioramento dell'esponente si traduce direttamente in un lower bound migliore: un bound $\log\xi\le O(k^{\alpha})$ con $\alpha<1/2$ darebbe $\chi\gtrsim n^{1/\alpha}$, e un bound polinomiale $\xi\le\poly(k)$ darebbe $\chi=2^{\Omega(n)}$. Gli autori pongono esplicitamente quest'ultima domanda \cite{SL26}.

\paragraph{Lower bound esponenziali noti in altri contesti.}
\begin{itemize}[leftmargin=1.5em]
\item \emph{Stati prodotto ``aritmeticamente generici''.} Lovitz e Steffan costruiscono stati prodotto espliciti (con fattori diversi) con $\chi\ge 2^n/(4n)$ ma $\chi_\delta=O(1)$ \cite{LS22}. Il meccanismo è aritmetico. Non si applica a $\Tst^{\otimes n}$, le cui ampiezze stanno nello stesso campo $\Q(\zeta_8)$ di quelle degli stati stabilizer (\cref{rem:arith}).
\item \emph{Fermioni.} Il rango gaussiano di $\ket M^{\otimes k}$, con $\ket M=(\ket{0000}+\ket{1111})/\sqrt2$, è almeno $\approx1.40^k$ \cite{RS26}. Il border rank approssimato è almeno $c_\delta(343/243)^k$ \cite{WWLL26}.
\item \emph{Bosoni.} Il rango in stati coerenti di $\ket1^{\otimes n}$ è $n^{\Omega(\log n)}$ \cite{CC26}, e il border rank è esattamente $2^n$ \cite{RS26}.
\end{itemize}
In tutti questi casi lo stato bersaglio ha una struttura (entanglement, o ampiezze trascendenti) che gli stati liberi non hanno. In più, gli stati liberi formano un'orbita continua con una struttura comune di annichilatori. La \cref{sec:barriers} spiega perché nessuna delle due cose è disponibile per $\Tst^{\otimes n}$.

\section{Un lower bound esponenziale condizionato}\label{sec:cond}

L'osservazione seguente è standard, nello spirito dei risultati di difficoltà della simulazione forte. La riportiamo per completezza, con le ipotesi esplicite.

\begin{hypothesis}[\#ETH non uniforme, versione per clausole]\label{hyp:nueth}
Esiste $\delta>0$ con questa proprietà: per ogni $M$ sufficientemente grande, ogni circuito booleano che conta le soluzioni delle formule 3-CNF con $M$ clausole ha taglia almeno $2^{\delta M}$.
\end{hypothesis}

\noindent La versione uniforme è \#ETH. Che la versione per variabili implichi quella per clausole segue da un lemma di sparsificazione per il conteggio \cite{DHMTW14}. La versione non uniforme è un'ipotesi più forte, ma usata in complessità fine-grained.

\begin{proposition}\label{prop:cond}
\leavevmode
\begin{enumerate}[(a),leftmargin=2em]
\item Sotto l'\cref{hyp:nueth} esiste $c>0$ tale che $\chi(\Tst^{\otimes t})\ge 2^{ct}$ per ogni $t$ sufficientemente grande.
\item (Incondizionato) Se $\chi(\Tst^{\otimes t})\le\poly(t)$, allora $\#\mathsf P\subseteq\mathsf{FP/poly}$. Di conseguenza $\mathsf P^{\#\mathsf P}=\mathsf{MA}$, e quindi $\mathsf{PH}\subseteq\mathsf{MA}$.
\end{enumerate}
\end{proposition}

\begin{proof}
\emph{Gadget.} Per uno stato a un qubit $\ket\psi=a\ket0+b\ket1$ vale
\[
T\ket\psi=\sqrt2\,(I\otimes\bra0)\,\mathrm{CNOT}\,(\ket\psi\otimes\Tst),
\]
con il CNOT controllato dal qubit dati. Infatti $\mathrm{CNOT}(\ket\psi\otimes\Tst)=\frac{1}{\sqrt2}\bigl[a\ket0(\ket0+\omega\ket1)+b\ket1(\ket1+\omega\ket0)\bigr]$ con $\omega=e^{i\pi/4}$, e la proiezione sul secondo qubit in $\ket0$ dà $(a\ket0+\omega b\ket1)/\sqrt2$.

Iterando il gadget, per un circuito Clifford$+T$ $U$ su $m$ qubit con $t$ gate $T$ esiste un Clifford $W$ tale che
\[
\bra{x}U\ket{0^m}=2^{t/2}\,\bra{x,0^t}\,W\,\bigl(\ket{0^m}\otimes\Tst^{\otimes t}\bigr).
\]

\emph{Advice.} Sia $\Tst^{\otimes t}=\sum_{j=1}^kc_j\ket{s_j}$. Allora
\[
\bra xU\ket{0^m}=2^{t/2}\sum_{j=1}^k c_j\,\bra{x,0^t}W\ket{0^m,s_j}.
\]
Ogni termine è un'ampiezza dello stato stabilizer $W\ket{0^m,s_j}$. Si calcola esattamente in $\Q(\zeta_8)$ in tempo $\poly(m+t)$, ad esempio con la CH-form \cite{BBCCGH19}. Possiamo scegliere gli $s_j$ linearmente indipendenti. Allora i $c_j$ sono l'unica soluzione di un sistema lineare a coefficienti in $\Q(\zeta_8)$, e per Cramer hanno una rappresentazione esatta di $\poly(k,t)$ bit.

La decomposizione dipende solo da $t$, quindi si può fornire come advice. Ogni ampiezza si calcola allora esattamente con un circuito di taglia $\poly(k,m,t)$.

\emph{Riduzione.} Sia $\varphi$ una 3-CNF con $N\le 3M$ variabili e $M$ clausole. Esiste un circuito reversibile $V$ con $O(M)$ gate di Toffoli che scrive $\varphi(x)$ in un registro ausiliario; ogni Toffoli si realizza con $7$ gate $T$ più Clifford. Il circuito $U=H^{\otimes N}V^\dagger Z_{\rm out}VH^{\otimes N}$ ha quindi $t=O(M)$ gate $T$, e
\[
\bra0U\ket0=2^{-N}\sum_x(-1)^{\varphi(x)}=1-2^{1-N}\,\#\varphi .
\]

(a) Se valesse $\chi(\Tst^{\otimes t(M)})\le 2^{\eps M}$ con $\eps$ abbastanza piccolo rispetto a $\delta$, si conterebbero le soluzioni di $\varphi$ con circuiti di taglia $2^{O(\eps M)}\poly(M)<2^{\delta M}$, contro l'ipotesi. Quindi $\chi(\Tst^{\otimes t(M)})\ge 2^{\Omega(M)}$ con $t(M)=O(M)$. La monotonia \ref{R3} estende il bound a ogni $t$.

(b) Con $k=\poly(t)$ si ottiene $\#3\mathsf{SAT}\in\mathsf{FP/poly}$. Poiché \#3SAT è $\#\mathsf P$-completo, $\#\mathsf P\subseteq\mathsf{FP/poly}$. Il resto è l'argomento standard: il permanente ha un protocollo interattivo con prover in $\mathsf P^{\#\mathsf P}$ \cite{LFKN92}, e se il permanente ha circuiti polinomiali Merlin può inviarli ad Arthur. Si ottiene $\mathsf P^{\#\mathsf P}\subseteq\mathsf{MA}$, e $\mathsf{PH}\subseteq\mathsf{MA}$ per il teorema di Toda.
\end{proof}

\begin{remark}
L'argomento non dice nulla sul rango approssimato con $\delta$ costante, perché il fattore $2^{t/2}$ del gadget amplifica l'errore. Questo è coerente con $\chi_\delta\le 1+\xi/\delta^2=O(2^{0.2284t})$.

Inoltre l'implicazione va in un solo verso. Un lower bound incondizionato su $\chi$ non darebbe separazioni di classi di complessità. Il problema è un lower bound per un modello di calcolo algebrico ristretto, cioè le combinazioni lineari di funzioni stabilizer, su una funzione esplicita: $x\mapsto e^{i\pi|x|/4}$.
\end{remark}

\section{Lower bound esponenziali incondizionati in modelli ristretti}\label{sec:restricted}

In tutta la sezione $\phi\in\C^2$ è uno stato normalizzato con $\phi\notin\Stab_1$. I risultati valgono quindi per $\Tst$, per $\Hst$ e per lo stato ``faccia''. Indichiamo con $d^\perp\in\Stab_1$ lo stato ortogonale a $d\in\Stab_1$.

\subsection{Un lemma di contrazione}

\paragraph{Fatto (marginali a un qubit).} Sia $s\in\Stab_n$ e sia $i$ un qubit. Lo stato ridotto è $\rho_i=\tfrac12\bigl(I+\sum_{P\in\{X,Y,Z\}}\langle P_i\rangle_s P\bigr)$, con $\langle P_i\rangle_s\in\{0,\pm1\}$. Al più uno dei tre valori è non nullo, perché $X_i,Y_i,Z_i$ anticommutano a coppie. Quindi:
\begin{itemize}[leftmargin=1.5em]
\item o $\rho_i=\ket d\bra d$ con $d\in\Stab_1$. Diciamo allora che il qubit $i$ è \emph{libero} in $s$, nello stato $d$. In questo caso $\ket s=\ket d_i\otimes\ket{s'}$ con $s'\in\Stab_{n-1}$;
\item oppure $\rho_i=I/2$. Diciamo allora che il qubit $i$ è \emph{entangled} in $s$.
\end{itemize}
Per $u\in\Stab_1$, la contrazione $(\bra u_i\otimes I)\ket s$ ha norma al quadrato $\bra u\rho_i\ket u\in\{0,\tfrac12,1\}$. Quando non è nulla, è proporzionale a uno stato stabilizer su $n-1$ qubit: lo stato post-misura della misura di Pauli che ha $u$ come autostato. Si annulla se e solo se il qubit $i$ è libero nello stato $u^\perp$.

\begin{lemma}[contrazione]\label{lem:contr}
Sia $\phi^{\otimes n}=\sum_{j=1}^kc_j\ket{s_j}$ con $s_j\in\Stab_n$. Fissiamo un qubit $i$ e $d\in\Stab_1$, e sia $K=\{j:\ \text{il qubit $i$ è libero in $s_j$ nello stato $d$}\}$. Allora $\phi^{\otimes(n-1)}$ ha una decomposizione in al più $k-|K|$ stati stabilizer. I nuovi termini sono le contrazioni dei vecchi: ogni qubit diverso da $i$ che era libero in $s_j$ resta libero.
\end{lemma}

\begin{proof}
Si applica $\bra{d^\perp}_i\otimes I$ ad ambo i membri. A sinistra si ottiene $\braket{d^\perp}{\phi}\,\phi^{\otimes(n-1)}$, e $\braket{d^\perp}{\phi}\neq0$ perché $\phi\not\propto d$. A destra i termini con $j\in K$ si annullano, mentre gli altri diventano multipli non nulli di stati stabilizer, per il Fatto. Si divide infine per $\braket{d^\perp}\phi$.

Quanto ai qubit liberi: se $\ket{s_j}=\bigotimes_{q\in U_j}\ket{d_q}\otimes\ket{r}$, dove $U_j$ è l'insieme dei qubit liberi, la contrazione agisce solo sul fattore che contiene il qubit $i$.
\end{proof}

\subsection{Decomposizioni in stati prodotto}

Sia $\ChiP(\psi)$ il minimo numero di termini in una decomposizione di $\psi$ in stati stabilizer prodotto, cioè in prodotti tensoriali di elementi di $\Stab_1$.

\begin{theorem}\label{thm:prod}
Sia $\phi\notin\Stab_1$ e sia $r_n=\ChiP(\phi^{\otimes n})$. Allora:
\begin{enumerate}[(i),leftmargin=2em]
\item $r_1=2$ e $r_n\ge\lceil\tfrac65\,r_{n-1}\rceil$;
\item $r_n=n+1$ per $1\le n\le 5$;
\item $6\,(6/5)^{\,n-5}\le r_n\le 6^{\lceil n/5\rceil}$ per $n\ge5$. In particolare $1.2\le\liminf r_n^{1/n}\le\limsup r_n^{1/n}\le 6^{1/5}\approx1.431$.
\end{enumerate}
\end{theorem}

\begin{proof}
(i) In una decomposizione prodotto ogni qubit di ogni termine è libero, e la contrazione conserva la forma prodotto. Sia $n_d$ il numero di termini il cui primo fattore è $d$. Il \cref{lem:contr} con $i=1$ dà $r_{n-1}\le r_n-n_d$ per ciascuno dei sei $d\in\Stab_1$. Sommando si ottiene $6r_{n-1}\le 6r_n-\sum_dn_d=5r_n$. Infine $r_1=2$, perché $\phi\notin\Stab_1$ e due stati stabilizer distinti generano $\C^2$.

(ii) Dal punto (i) seguono i lower bound $r_2\ge3$, $r_3\ge4$, $r_4\ge5$, $r_5\ge6$. Per l'upper bound: $\phi^{\otimes n}$ sta nel sottospazio simmetrico $\mathrm{Sym}^n(\C^2)$, che ha dimensione $n+1$. I sei vettori $d^{\otimes n}$, $d\in\Stab_1$, corrispondono a sei punti distinti della curva razionale normale. Quindi $\min(6,n+1)$ qualsiasi di essi sono linearmente indipendenti. Per $n\le5$, $n+1$ di essi generano $\mathrm{Sym}^n$, e $r_n\le n+1$.

(iii) Il lower bound si ottiene iterando (i) a partire da $r_5=6$. L'upper bound segue dalla submoltiplicatività $r_{a+b}\le r_ar_b$ e da $r_m\le6$ per $m\le5$.
\end{proof}

\begin{remark}
Le ricerche esaustive per $\phi=\Hst$ confermano $r_2=3$ e $r_3=4$ (\cref{app:num}). Il modello prodotto perde già a $n=2$: con termini entangled si ha
\[
\Hst^{\otimes2}=\tfrac1{\sqrt2}\bigl(\ket{\Phi^+}+\mathrm{CZ}\ket{++}\bigr),
\]
e i due stati sono ortogonali. Quindi $\chi(\Hst^{\otimes2})=2<3=r_2$.
\end{remark}

\subsection{Termini con molti qubit liberi}

\begin{theorem}\label{thm:free}
Sia $\phi^{\otimes n}=\sum_{j=1}^kc_j\ket{s_j}$, e supponiamo che ogni $s_j$ abbia almeno $m$ qubit liberi. Allora
\[
k\ \ge\ \exp\!\Bigl(\frac{m(m+1)}{12\,n}\Bigr).
\]
\end{theorem}

\begin{proof}
Costruiamo una successione di decomposizioni $\mathcal D_0,\dots,\mathcal D_m$. La decomposizione $\mathcal D_\sigma$ è di $\phi^{\otimes(n-\sigma)}$, ha $k_\sigma$ termini, e ogni suo termine ha almeno $m-\sigma$ qubit liberi. Si parte da $\mathcal D_0$, la decomposizione data, con $k_0=k$.

Al passo $\sigma$ si scelgono a caso un qubit $i$, uniforme fra gli $n-\sigma$ rimasti, e uno stato $d$, uniforme in $\Stab_1$. Un termine con $u\ge m-\sigma$ qubit liberi viene eliminato dal \cref{lem:contr} con probabilità $\frac{u}{6(n-\sigma)}\ge\frac{m-\sigma}{6(n-\sigma)}$. Esiste quindi una scelta di $(i,d)$ che elimina almeno la frazione attesa di termini, e dà
\[
k_{\sigma+1}\le k_\sigma\Bigl(1-\frac{m-\sigma}{6(n-\sigma)}\Bigr).
\]
I termini superstiti hanno almeno $m-\sigma-1$ qubit liberi, per il \cref{lem:contr}.

Dopo $m$ passi si ha $k_m\ge1$, quindi
\[
1\le k\prod_{\sigma=0}^{m-1}\Bigl(1-\frac{m-\sigma}{6(n-\sigma)}\Bigr)\le k\exp\Bigl(-\sum_{\sigma=0}^{m-1}\frac{m-\sigma}{6n}\Bigr)=k\,e^{-m(m+1)/(12n)}.\qedhere
\]
\end{proof}

\begin{corollary}\label{cor:free}
Ogni decomposizione di $\phi^{\otimes n}$ in $k$ stati stabilizer contiene un termine con meno di $\sqrt{12\,n\ln k}$ qubit liberi. In particolare:
\begin{itemize}[leftmargin=1.5em]
\item se $k=n^{O(1)}$, almeno un termine ha $n-O(\sqrt{n\log n})$ qubit entangled;
\item se ogni termine ha almeno $\alpha n$ qubit liberi, allora $k\ge e^{\alpha^2n/12}$.
\end{itemize}
\end{corollary}

\subsection{Decomposizioni ben condizionate}

\begin{proposition}\label{prop:cond-number}
Sia $\phi^{\otimes n}=\sum_{j=1}^kc_j\ket{s_j}$ con gli $s_j$ linearmente indipendenti, e sia $\Gram_{ij}=\braket{s_i}{s_j}$ la loro matrice di Gram. Allora
\[
k\ \ge\ \lambda_{\min}(\Gram)\,\xi(\phi)^n .
\]
Per $\phi=\Hst$ questo dà $k\ge\lambda_{\min}(\Gram)\,(4-2\sqrt2)^n\approx\lambda_{\min}(\Gram)\,2^{0.2284n}$.
\end{proposition}

\begin{proof}
Si ha $1=c^\dagger\Gram c\ge\lambda_{\min}\|c\|_2^2$, quindi $\xi(\phi^{\otimes n})\le\|c\|_1^2\le k\|c\|_2^2\le k/\lambda_{\min}$. Si conclude con la moltiplicatività $\xi(\phi^{\otimes n})=\xi(\phi)^n$ \cite{BBCCGH19}.
\end{proof}

\noindent Per esempio, una decomposizione in stati ortogonali richiede almeno $\xi(\phi)^n$ termini. Più in generale, una decomposizione con $\poly(n)$ termini deve avere $\lambda_{\min}(\Gram)\le\poly(n)\cos^{2n}(\pi/8)$ e $\|c\|_2^2\ge\xi/k$: famiglie di stati stabilizer esponenzialmente vicine alla dipendenza lineare, e cancellazioni esponenziali fra i coefficienti.

\begin{remark}[blocchi di taglia fissa]\label{rem:blocks}
Supponiamo che tutti i termini siano prodotti tensoriali di stati stabilizer sui blocchi di una \emph{stessa} partizione di $[n]$ in blocchi di taglia $\le b$. Lo stesso argomento, con funzionali stabilizer su un blocco intero, dà $k\ge(1-p_b)^{-\lceil n/b\rceil}$ con $p_b\ge1/|\Stab_b|>0$.

Per vederlo, si contrae un blocco $B$ alla volta con $\omega$ uniforme in $\Omega_B=\{\omega\in\Stab_{|B|}:\braket{\omega}{\phi^{\otimes|B|}}\neq0\}$. Un termine il cui fattore su $B$ è $\sigma$ viene eliminato se $\omega\perp\sigma$. Gli stati della base stabilizer di $\sigma$ diversi da $\sigma$ generano $\sigma^\perp$. Poiché $\phi^{\otimes|B|}\not\propto\sigma$, almeno uno di essi sta in $\Omega_B$, e quindi la probabilità di eliminazione è almeno $1/|\Omega_B|\ge1/|\Stab_b|$.

Il bound decade però rapidamente con $b$, e non copre il caso in cui la partizione cambia da termine a termine.
\end{remark}

\subsection{Perché il metodo si ferma}\label{sec:barriers}

Le osservazioni seguenti non sono teoremi di impossibilità. Spiegano però perché le tecniche ``semplici'' non arrivano a $\Tst^{\otimes n}$ con termini arbitrari.

\begin{enumerate}[label=(B\arabic*),leftmargin=2.5em]
\item\label{B1} \textbf{I qubit entangled non si eliminano.} Se il qubit $i$ è entangled in $s$, allora $\rho_i=I/2$ e $(\bra u\otimes I)\ket s\neq0$ per ogni $u$.

Più in generale, sia $A$ un insieme di qubit e $\omega$ uno stato su $A$. Allora $(\bra\omega_A\otimes I)\ket s=0$ se e solo se $\omega\perp\operatorname{range}\rho_A(s)$. Per uno stato stabilizer, $\rho_A$ è proporzionale al proiettore su un codice stabilizer di dimensione $2^{S_A}$, dove $S_A$ è l'entropia di entanglement in bit. Se $A$ è massimamente entangled con il complementare, nessun funzionale lo elimina, e questo accade tipicamente per $|A|<n/2$.

Già a $n=2$ la decomposizione ottima usa stati entangled (Bell e $\mathrm{CZ}\ket{++}$).

\item\label{B2} \textbf{Flattening.} Le tecniche che funzionano per fermioni e bosoni sono mappe lineari $\Phi$ (matrici di monomi di Majorana o di operatori di annichilazione applicati allo stato) con due proprietà: il rango è subadditivo, ed è piccolo su \emph{tutti} gli stati liberi simultaneamente.

La seconda proprietà viene da una struttura comune di annichilatori. Gli stati gaussiani sono vuoti di modi ruotati; gli stati coerenti sono autovettori simultanei degli operatori di annichilazione \cite{RS26,WWLL26,CC26}. Gli stati stabilizer invece formano un insieme finito: ciascuno è autovettore del \emph{proprio} gruppo stabilizzatore, ma non esiste una famiglia comune.

Inoltre $\phi^{\otimes n}$ è uno stato prodotto. I flattening basati su bipartizioni (rango di Schmidt) valgono $1$ su $\phi^{\otimes n}$. Gli spazi $\operatorname{span}\{O\psi: O\ \text{operatore locale di grado}\le p\}$ sono minimi proprio sugli stati prodotto. In entrambi i casi la disuguaglianza va nel verso sbagliato.

\item\label{B3} \textbf{Proiezione simmetrica.} Lo stato $\phi^{\otimes n}$ vive in $\mathrm{Sym}^n(\C^2)$, che ha dimensione $n+1$. Un argomento che usa solo le proiezioni $\Pi_{\rm sym}s_j$ dei termini dà al più $n+1$.

\item\label{B4} \textbf{Aritmetica.} Si veda l'osservazione seguente.
\end{enumerate}

\begin{remark}[il meccanismo aritmetico]\label{rem:arith}
Sia $K=\Q(\zeta_8)$. Ogni stato stabilizer ha componenti in $K$, perché $1/\sqrt2\in K$. Sia $\ket\psi=\bigotimes_{q=1}^n(\ket0+\alpha_q\ket1)$ con $\alpha_1,\dots,\alpha_n$ algebricamente indipendenti su $K$.

Supponiamo che $\psi\in V=\operatorname{span}_\C\{s_1,\dots,s_k\}$. Sia $\pi$ una mappa $K$-lineare con nucleo $V_K=\operatorname{span}_K\{s_j\}$. Allora $0=\pi(\psi)=\sum_x\alpha^x\pi(e_x)$, dove $\alpha^x=\prod_q\alpha_q^{x_q}$. I monomi multilineari $\alpha^x$ sono linearmente indipendenti su $K$, quindi $\pi(e_x)=0$ per ogni $x$. Dunque ogni $e_x$ sta in $V$, e $\chi(\psi)=2^n$.

Se invece i fattori sono tutti uguali, $\alpha_q=\alpha$ trascendente, si ottiene solo che i vettori di Dicke stanno in $V$, cioè $\chi\ge n+1$. Per $\Hst$ si ha $\alpha=\tan\frac\pi8=\sqrt2-1\in K$, e non si ottiene alcun vincolo. Versioni esplicite (algebriche) del primo caso si trovano in \cite{LS22}.

In sintesi, per $\Tst^{\otimes n}$ la via aritmetica è chiusa: lo stato è ``definito sullo stesso campo'' degli stati stabilizer, ed è invariante per un gruppo di simmetrie grande (permutazioni e $H^{\otimes n}$).
\end{remark}


\section{Riformulazione: entropia stabilizer e norme di Gowers}\label{sec:sre}

Per $\psi\in\C^{2^n}$ sia $f_\psi:\Fb_2^n\to\C$, $f_\psi(x)=2^{n/2}\braket{x}{\psi}$, così che $\E_x|f_\psi(x)|^2=1$. Indichiamo con $\mathcal C$ la coniugazione complessa. La norma di Gowers $U^3$ è
\[
\|f\|_{U^3}^8=\E_{x,h_1,h_2,h_3\in\Fb_2^n}\ \prod_{\omega\in\{0,1\}^3}\mathcal C^{|\omega|}f\bigl(x+\omega_1h_1+\omega_2h_2+\omega_3h_3\bigr).
\]

\begin{proposition}[SRE $=$ Gowers $U^3$]\label{prop:gowers}
Per ogni stato normalizzato $\psi$,
\[
2^{-M_2(\psi)}=\|f_\psi\|_{U^3}^8,\qquad\text{cioè}\qquad M_2(\psi)=-8\log_2\|f_\psi\|_{U^3}.
\]
\end{proposition}

\noindent La quantità $\|f_\psi\|_{U^3}$ è la ``norma di Gowers di uno stato quantistico'' di \cite{AD24}. Lì il legame con il quarto momento dello spettro di Pauli è implicito; qui lo rendiamo esplicito.

\begin{proof}
A meno di una fase, $\bra\psi X^aZ^b\ket\psi=\hat g_a(b)$, dove $g_a(x)=\overline{\psi(x+a)}\,\psi(x)$ e $\hat g(b)=\sum_xg(x)(-1)^{b\cdot x}$. Per ogni $g$ vale l'identità
\[
\sum_b|\hat g(b)|^4=2^n\sum_{x,h_1,h_2}g(x)\,\overline{g(x+h_1)}\,\overline{g(x+h_2)}\,g(x+h_1+h_2).
\]
Inoltre $g_a=2^{-n}\Delta_af$, con $f=f_\psi$ e $\Delta_af(x)=f(x)\overline{f(x+a)}$. Sostituendo,
\[
2^{-n}\sum_{a,b}|\hat g_a(b)|^4=\E_a\|\Delta_af\|_{U^2}^4=\|f\|_{U^3}^8 .
\]
Il primo membro è $2^{-n}\sum_P\bra\psi P\ket\psi^4=2^{-M_2(\psi)}$.
\end{proof}

\begin{proposition}\label{prop:M2F}
Per ogni stato normalizzato $\psi$,
\[
\tfrac14M_2(\psi)\ \le\ \log_2\frac1{F(\psi)}\ \le\ \log_2\xi(\psi).
\]
\end{proposition}

\begin{proof}
Sia $s$ uno stato stabilizer che realizza $F$, e sia $S$ il suo gruppo stabilizzatore con segni, formato da $2^n$ operatori di Pauli con $g\ket s=\ket s$. Allora $\ket s\bra s=2^{-n}\sum_{g\in S}g$, e quindi $F=2^{-n}\sum_{g\in S}\bra\psi g\ket\psi$. Poiché $x\mapsto x^4$ è convessa, la disuguaglianza di Jensen dà
\[
2^{-n}\sum_{P}\bra\psi P\ket\psi^4\ \ge\ 2^{-n}\sum_{g\in S}\bra\psi g\ket\psi^4\ \ge\ \Bigl(2^{-n}\sum_{g\in S}\bra\psi g\ket\psi\Bigr)^4=F^4.
\]
La seconda disuguaglianza dell'enunciato è \ref{R2}.
\end{proof}

\noindent Vale anche una relazione inversa polinomiale. Il teorema inverso per la norma di Gowers $U^3$ degli stati quantistici \cite{AD24}, che passa per la congettura di Marton (PFR) su $\Fb_2^n$ \cite{GGMT23}, dà $F(\psi)\ge\poly\bigl(2^{-M_2(\psi)}\bigr)$. Quindi $\log_2(1/F)$ e $M_2$ coincidono a meno di costanti moltiplicative e additive. Le costanti in \cite{AD24} sono grandi; qui conta solo che la relazione sia polinomiale.

\subsection{Tre congetture equivalenti a meno di costanti}

\begin{conjecture}\label{conj:main}
Esistono costanti $C,D$ tali che, per ogni $n$ e ogni stato puro $\psi$ su $n$ qubit:
\begin{enumerate}[(A),leftmargin=2em]
\item\label{cA} (SRE--rango) $M_2(\psi)\le C\log_2\chi(\psi)+D$;
\item\label{cB} (fidelity--rango) $\log_2(1/F(\psi))\le C\log_2\chi(\psi)+D$;
\item\label{cC} (extent--rango, domanda aperta in \cite{SL26}) $\log_2\xi(\psi)\le C\log_2\chi(\psi)+D$.
\end{enumerate}
\end{conjecture}

\noindent Per la \cref{prop:M2F} si ha \ref{cC}$\Rightarrow$\ref{cB}$\Rightarrow$\ref{cA}, e per il teorema inverso di \cite{AD24} vale anche \ref{cA}$\Rightarrow$\ref{cB}, con costanti diverse.

\begin{theorem}\label{thm:conj-exp}
Ciascuna delle tre congetture implica $\chi(\phi^{\otimes n})=2^{\Omega(n)}$ per ogni $\phi\notin\Stab_1$. Per $\Hst$, la \ref{cA} dà $\chi(\Hst^{\otimes n})\ge2^{(n\log_2(4/3)-D)/C}$ e la \ref{cB} dà $\chi(\Hst^{\otimes n})\ge2^{(0.2284\,n-D)/C}$.
\end{theorem}

\begin{proof}
$M_2$ è additiva e strettamente positiva su $\phi\notin\Stab_1$. Fidelity ed extent sono moltiplicativi sui prodotti di stati a un qubit \cite{BBCCGH19}. Si usano poi i valori \eqref{eq:values}.
\end{proof}

\paragraph{Cosa si sa.} Il risultato di Sinha--Lovitz, insieme alla \cref{prop:M2F}, dà $M_2(\psi)\le O\bigl(\sqrt{k\log k}\bigr)$ per stati di rango $k$. Kalra--Sinha danno $\log_2(1/F)\le O(k\log k)$ \cite{KS25}. Tutto il gap fra il lower bound quadratico e quello esponenziale sta nella differenza fra $\sqrt k$ e $\log k$.

\paragraph{Vincoli sulle costanti.}
\begin{itemize}[leftmargin=1.5em]
\item Per somme di $k$ stati stabilizer quasi ortogonali con pesi uguali, i contributi diagonali danno $2^{-M_2}\approx\sum_j|c_j|^8=k^{-3}$, cioè $M_2\approx3\log_2k$. Numericamente, a $n=12$, si trova $M_2=2.97,\ 4.74,\ 5.93$ per $k=2,3,4$ (\cref{app:num}). Quindi serve $C\ge3$ in \ref{cA}, e $C\ge1$ in \ref{cB}, perché in quel caso $F\approx1/k$.
\item \emph{Forma forte.} Si potrebbe congetturare $F(\psi)\,\chi(\psi)\ge1$, cioè \ref{cB} con $C=1$ e $D=0$. Darebbe $\chi(\Hst^{\otimes n})\ge(4-2\sqrt2)^n\approx1.1716^n$, che coincide con lo scaling del rango approssimato. I dati sono vicini al limite:
\begin{itemize}
\item per $\Hst$: $F^m\chi=1.71,\ 1.46,\ 2.12$ per $m=1,2,4$;
\item per lo stato faccia, con $F=\tfrac12(1+1/\sqrt3)\approx0.789$ e i valori di $\chi$ di \cite{LR26}: $F^m\chi=1.58,\ 1.24,\ 1.47,\ 1.16$ per $m=1,\dots,4$;
\item per lo stato $W_n$: $F(W_n)=9/(4n)$, mentre $\chi(W_n)\le\lceil n/2\rceil$ (somme di coppie), quindi $F\chi\le9/8+O(1/n)$.
\end{itemize}
Una decomposizione delle potenze dello stato faccia con esponente minore di $\log_2(1/F)\approx0.342$ per copia confuterebbe la forma forte, ma non \ref{cB}.
\end{itemize}

\subsection{Formulazione additivo-combinatoria}

\begin{conjecture}[forma Gowers]\label{conj:gowers}
Sia $f=\sum_{j=1}^kc_jg_j$, dove ogni $g_j$ è una funzione stabilizer normalizzata in $L^2$: una fase quadratica, eventualmente non classica a valori in $\Z_4$, ristretta a un sottospazio affine di $\Fb_2^n$. Se $\E|f|^2=1$, allora $\|f\|_{U^3}\ge k^{-O(1)}$.
\end{conjecture}

\noindent Per la \cref{prop:gowers} questa è esattamente la \ref{cA}. Alcune osservazioni.
\begin{itemize}[leftmargin=1.5em]
\item Il bound noto è $\|f\|_{U^3}^8\ge F^4\ge\xi^{-4}\ge2^{-O(\sqrt{k\log k})}$ \cite{SL26}.
\item L'analogo per $U^2$ con i caratteri è banale: se $f=\sum_jc_j\chi_j$ con caratteri ortonormali, allora $\|f\|_{U^2}^4=\sum_j|c_j|^4\ge1/k$.
\item La difficoltà sta in due punti. Primo, $f$ non è limitata ($|f|$ arriva a $2^{n/2}$), mentre i teoremi inversi standard sono per funzioni limitate. Secondo, le cancellazioni fra i $g_j$ devono essere esponenziali (\cref{prop:cond-number}).
\item Una risposta positiva dà $\chi(\Tst^{\otimes n})=2^{\Omega(n)}$. Un controesempio, cioè somme di $k$ funzioni stabilizer con norma $U^3$ pari a $k^{-\omega(1)}$, mostrerebbe che nessuna disuguaglianza generale fra rango e $M_2$ (o fidelity, o extent) può dare un bound esponenziale.
\end{itemize}

\section{Problemi aperti e piano di lavoro}\label{sec:open}

Definiamo le funzioni
\[
f(k)=\sup_{n}\ \sup_{\psi:\ \chi(\psi)\le k}M_2(\psi),\qquad
g(k)=\sup_{n}\ \sup_{\psi:\ \chi(\psi)\le k}\log_2\frac1{F(\psi)} .
\]
Esse sono finite e indipendenti da $n$ \cite{SL26}. Per quanto visto,
\[
3\log_2k-O(1)\ \lesssim\ f(k)\ \le\ 4g(k)\ \le\ O\bigl(\sqrt{k\log k}\bigr),
\qquad \log_2k-O(1)\ \lesssim\ g(k).
\]
Il lower bound su $f$ è numerico/euristico (\cref{app:num}). La \cref{conj:main} dice che $f(k)=O(\log k)$; per il \cref{thm:conj-exp}, questo basta per l'esponenziale.

\begin{enumerate}[label=\textbf{P\arabic*.},leftmargin=2.2em]
\item \textbf{Crescita di $f$ e $g$.} Dimostrare $f(k)=O(\log k)$, oppure trovare famiglie con $f(k)=\omega(\log k)$. Un primo passo concreto è dimostrare la \ref{cB} per decomposizioni con Gram a blocchi, cioè unioni di poche famiglie ben condizionate. L'idea è rendere quantitativo il passo di ``fusione dei codespace'' di \cite{SL26}, che è dove si perde $\sqrt k$.

\item \textbf{Ricerca di controesempi.} Massimizzare $M_2$, o minimizzare $F\cdot k$, su stati di rango $k\le 6$ per $n\le 5$: ricerca esaustiva su coppie e terne, ottimizzazione continua sui coefficienti. Poi testare famiglie strutturate: stati $W$ e di Dicke, sovrapposizioni di stati cat, stati ipergrafo di rango basso. Un controesempio superpolinomiale chiuderebbe la strada della \cref{sec:sre}, che è comunque un risultato.

\item \textbf{Valori esatti.} Calcolare $\chi(\Tst^{\otimes m})$ e $\chi(\ket{F}^{\otimes m})$ (stato faccia) per $m=5,\dots,8$. Il \cref{lem:contr} dà vincoli ricorsivi utilizzabili per potare la ricerca, e le simmetrie riducono lo spazio: permutazioni dei qubit, e Clifford a un qubit che fissano $\phi$ applicati a tutti i qubit. Questi dati misurano l'esponente e mettono alla prova la forma forte $F\chi\ge1$.

\item \textbf{Modelli ristretti più ricchi.}
\begin{itemize}
\item prodotti a blocchi con partizioni che variano da termine a termine, dove si incontra l'ostacolo \ref{B1};
\item termini con entanglement limitato, ad esempio $S_A\le s$ per ogni $|A|\le a$;
\item decomposizioni reali o CSS per $\Hst$, che ha ampiezze reali;
\item coefficienti limitati: da $\xi\le k^2\|c\|_\infty^2$ si ha banalmente $k\ge\sqrt{\xi}/\|c\|_\infty$.
\end{itemize}

\item \textbf{Modello prodotto.} Determinare $\lim r_n^{1/n}\in[1.2,\,1.431]$; per esempio, $r_6\in[8,12]$.

\item \textbf{Gate.} Stabilire se l'unitarietà dei termini aiuta nella decomposizione del gate, $T^{\otimes n}=\sum_jc_jC_j$. La versione operatoriale della \ref{cC} riguarda $\Lambda_\Cl(A)^2$ contro $\chi_\Cl(A)$ \cite{SL26}.
\end{enumerate}

\paragraph{Priorità suggerita.} Conviene partire da P2 e P3, perché sono calcoli finiti che possono cambiare il quadro: o rafforzano la \cref{conj:main}, o la confutano. P1 è il problema vero, e P4--P5 sono riscaldamenti con risultati pubblicabili in forma breve.

\appendix
\section{Verifiche numeriche}\label{app:num}

Tutti i calcoli sono in Python/NumPy. Gli stati stabilizer per $n=3$ (1080 stati) e $n=4$ (36720 stati) sono generati per visita in ampiezza a partire da $\ket{0^n}$, applicando $H$, $S$ e CNOT e identificando gli stati che differiscono per una fase globale.

\begin{enumerate}[leftmargin=1.5em]
\item \textbf{Decomposizioni esplicite.} Verificati:
\begin{itemize}
\item $\Hst^{\otimes2}\in\operatorname{span}\{\ket{00},\ket{11},\ket{++}\}$;
\item $\Hst^{\otimes5}\in\operatorname{span}\{d^{\otimes5}:d\in\Stab_1\}$;
\item $\Hst^{\otimes2}=(\ket{\Phi^+}+\mathrm{CZ}\ket{++})/\sqrt2$, con $\braket{\Phi^+}{\mathrm{CZ}\,{+}{+}}=0$.
\end{itemize}

\item \textbf{Rango prodotto (\cref{thm:prod}).} La ricerca esaustiva sui $36$ stati prodotto a $n=2$ dà $r_2=3$. A $n=3$, nessuna delle $\binom{216}{3}\approx1.7\cdot10^6$ terne di stati prodotto genera $\Hst^{\otimes3}$, quindi $r_3=4$. Entrambi i valori coincidono con quelli dedotti dalla ricorsione.

\item \textbf{Identità SRE--Gowers (\cref{prop:gowers}).} Su stati casuali a $3$ qubit, $2^{-M_2}$ e $\|f_\psi\|_{U^3}^8$ coincidono a $10^{-6}$. Per esempio, $0.396849$ contro $0.396849$.

\item \textbf{\cref{prop:M2F}.} Su $200$ stati casuali a $3$ qubit, con $F$ calcolata esattamente sui $1080$ stati stabilizer, il massimo di $M_2/\log_2(1/F)$ è $3.34$, sotto il bound $4$. Per $\Hst^{\otimes3}$ si trova $M_2=1.2451=3\log_2\frac43$ e $F=0.6219=\cos^6\frac\pi8$.

\item \textbf{Stato $W$.} $F(W_4)=0.5625=9/16$, calcolata esattamente sui $36720$ stati stabilizer a $4$ qubit.

\item \textbf{SRE di somme di $k$ stati stabilizer casuali} ($n=12$). Gli stati sono generati da circuiti Clifford casuali di profondità $400$, sommati con fasi casuali e normalizzati; $4$ campioni per ogni $k$. $M_2$ è calcolata esattamente con la trasformata di Walsh--Hadamard. Il valore tipico per uno stato di Haar, cioè la saturazione, è $\approx\log_2\frac{2^n+3}{4}\approx10$.
\begin{center}
\begin{tabular}{@{}rccccccc@{}}
\toprule
$k$ & 1 & 2 & 3 & 4 & 6 & 8\\
\midrule
$M_2$ & $0.00$ & $2.97\pm0.07$ & $4.74\pm0.05$ & $5.93\pm0.06$ & $7.46\pm0.06$ & $8.39\pm0.05$\\
$3\log_2k$ & $0.00$ & $3.00$ & $4.75$ & $6.00$ & $7.75$ & $9.00$\\
\bottomrule
\end{tabular}
\end{center}
Le deviazioni per $k\ge6$ sono dovute alla vicinanza della saturazione.
\end{enumerate}

\begin{thebibliography}{99}\small

\bibitem{BSS16} S.~Bravyi, G.~Smith, J.~A.~Smolin, \emph{Trading classical and quantum computational resources}, Phys.\ Rev.\ X \textbf{6}, 021043 (2016), arXiv:1506.01396.

\bibitem{BG16} S.~Bravyi, D.~Gosset, \emph{Improved classical simulation of quantum circuits dominated by Clifford gates}, Phys.\ Rev.\ Lett.\ \textbf{116}, 250501 (2016), arXiv:1601.07601.

\bibitem{BBCCGH19} S.~Bravyi, D.~Browne, P.~Calpin, E.~Campbell, D.~Gosset, M.~Howard, \emph{Simulation of quantum circuits by low-rank stabilizer decompositions}, Quantum \textbf{3}, 181 (2019), arXiv:1808.00128.

\bibitem{QPG21} H.~Qassim, H.~Pashayan, D.~Gosset, \emph{Improved upper bounds on the stabilizer rank of magic states}, Quantum \textbf{5}, 606 (2021), arXiv:2106.07740.

\bibitem{PSV22} S.~Peleg, A.~Shpilka, B.~L.~Volk, \emph{Lower bounds on stabilizer rank}, Quantum \textbf{6}, 652 (2022), arXiv:2106.03214.

\bibitem{LS22} B.~Lovitz, V.~Steffan, \emph{New techniques for bounding stabilizer rank}, Quantum \textbf{6}, 692 (2022), arXiv:2110.07781.

\bibitem{Lab22} F.~Labib, \emph{Stabilizer rank and higher-order Fourier analysis}, Quantum \textbf{6}, 645 (2022), arXiv:2107.10551.

\bibitem{MT23} S.~Mehraban, M.~Tahmasbi, \emph{Quadratic lower bounds on the approximate stabilizer rank: a probabilistic approach}, arXiv:2305.10277 (2023).

\bibitem{KS25} A.~R.~Kalra, P.~Sinha, \emph{Stabilizer ranks, Barnes--Wall lattices and magic monotones}, Quantum \textbf{10}, 2179 (2026), arXiv:2503.04101.

\bibitem{LR26} F.~Labib, V.~Russo, \emph{Stabilizer rank bounds for magic-state orbits}, arXiv:2605.28586 (2026).

\bibitem{SL26} P.~Sinha, B.~Lovitz, \emph{Improved bounds on stabilizer extent and Clifford rank}, arXiv:2610.06819 (2026).

\bibitem{CC26} F.~Cottier, U.~Chabaud, \emph{Lower bounds on coherent state rank}, arXiv:2604.00766 (2026).

\bibitem{RS26} O.~Reardon-Smith, \emph{Exponential lower bounds on the fermionic Gaussian rank of magic states and the bosonic coherent state rank of Fock states}, arXiv:2610.03699 (2026).

\bibitem{WWLL26} F.~Wei, K.-W.~Wu, Z.~Liu, Z.-W.~Liu, \emph{Robust exponential lower bounds for fermionic and bosonic Gaussian ranks}, arXiv:2610.02172 (2026).

\bibitem{LOH22} L.~Leone, S.~F.~E.~Oliviero, A.~Hamma, \emph{Stabilizer Rényi entropy}, Phys.\ Rev.\ Lett.\ \textbf{128}, 050402 (2022), arXiv:2106.12587.

\bibitem{AD24} S.~Arunachalam, A.~Dutt, \emph{Polynomial-time tolerant testing stabilizer states}, arXiv:2408.06289 (2024).

\bibitem{GGMT23} W.~T.~Gowers, B.~Green, F.~Manners, T.~Tao, \emph{On a conjecture of Marton}, arXiv:2311.05762 (2023).

\bibitem{DHMTW14} H.~Dell, T.~Husfeldt, D.~Marx, N.~Taslaman, M.~Wahlén, \emph{Exponential time complexity of the permanent and the Tutte polynomial}, ACM Trans.\ Algorithms \textbf{10}, 21 (2014).

\bibitem{LFKN92} C.~Lund, L.~Fortnow, H.~Karloff, N.~Nisan, \emph{Algebraic methods for interactive proof systems}, J.\ ACM \textbf{39}, 859 (1992).

\end{thebibliography}


\end{document}
